% Answers to selected exercises, chapter 01.

\paragraph{Exercise 1.1}
An orthogonal projector is symmetric, so it has an orthonormal eigenbasis. From
$\Pi^2=\Pi$, every eigenvalue satisfies $\lambda^2=\lambda$, so it is $0$ or
$1$. The trace is the number of eigenvalues equal to $1$, which is the dimension
of the range, the rank. For the projector of \cref{ex:ch01:gproj} in $\R^3$ the
diagonal is $(\tfrac12,\tfrac12,1)$, with sum $2$, and the range is spanned by
$q_1$ and $q_2$.

\paragraph{Exercise 1.2}
$\Pi^2=\begin{pmatrix}1&1\\0&0\end{pmatrix}=\Pi$, so it is a projector. It is not
symmetric, so it is not an orthogonal projector. Its range is spanned by $e_1$
and its kernel by $(1,-1)\T$. The angle between them is $45^\circ$, since
$e_1\cdot(1,-1)\T/\sqrt2=1/\sqrt2$. An orthogonal projector would have a
$90^\circ$ angle.

\paragraph{Exercise 1.3}
$A-B=\begin{pmatrix}1&-1\\-1&-1\end{pmatrix}$ has determinant $-2<0$, so one
eigenvalue is positive and one negative. Neither $A\psd B$ nor $B\psd A$ holds.

\paragraph{Exercise 1.4}
Take $A=\begin{pmatrix}1&0\\0&0\end{pmatrix}$, $B=\begin{pmatrix}0&1\\0&0\end{pmatrix}$
and $C=\begin{pmatrix}0&0\\1&0\end{pmatrix}$. Then $BC=\diag(1,0)$, so
$\tr(ABC)=1$. And $CB=\diag(0,1)$, so $ACB=0$ and $\tr(ACB)=0$.

\paragraph{Exercise 1.5}
$\Sigma_x^{1/2}g=(2,-1)\T$, so $\tilde P=(2,-1)\T(2,-1)=\begin{pmatrix}4&-2\\-2&1\end{pmatrix}$.
It has rank one, with eigenvalue $4+1=5$ along $(2,-1)\T/\sqrt5$. And
$\tr(P\Sigma_x)=g\T\Sigma_xg=4+1=5$, equal to $\tr\tilde P$.

\paragraph{Exercise 1.6}
$\det(A-\mu B)=(1-2\mu)(-\mu)-\mu^2=\mu^2-\mu$, so $\mu=0$ or $\mu=1$. For
$\mu=0$, $Av=0$ gives $v=(0,1)\T$. For $\mu=1$, $(A-B)v=0$ with
$A-B=\begin{pmatrix}-1&-1\\-1&-1\end{pmatrix}$ gives $v=(1,-1)\T$. The two are
$B$-orthogonal, since $B(1,-1)\T=(1,0)\T$ and $(0,1)\cdot(1,0)=0$. The maximum
of $x\T Ax/x\T Bx$ is $1$, at $x=(1,-1)\T$, where the numerator is $1$ and the
denominator is $2-2+1=1$.

\paragraph{Exercise 1.7}
With $A=2$ and $B=(1,1)$, $S=\begin{pmatrix}2&1\\1&2\end{pmatrix}-\tfrac12\begin{pmatrix}1&1\\1&1\end{pmatrix}
=\begin{pmatrix}1.5&0.5\\0.5&1.5\end{pmatrix}$. Its trace is $3$ and its
determinant is $2$, so it is positive definite, and $M\pd0$ by
\cref{thm:ch01:schurpsd}. Then $\det M=2\cdot2=4$. The eigenvalues of $M$ are
$4,1,1$, whose product is also $4$.

\paragraph{Exercise 1.8}
A sign flip sends $e_j$ to $s_je_j$, and a permutation sends that to $s_je_{\pi(j)}$,
a signed coordinate vector. All of its energy stays on one coordinate. For the
randomized Hadamard matrix, $D_se_j=s_je_j$ and $H_d(s_je_j)$ is $s_j$ times the
$j$-th column of $H_d$, whose entries are all $\pm1$. Dividing by $\sqrt d$ gives
coordinates of magnitude $1/\sqrt d$.

\paragraph{Exercise 1.9}
$\det(X+H)=8-0.01=7.99$, so the exact change is $\log(7.99/8)\approx-0.00125078$.
The first-order term is $\tr(X^{-1}H)=0$, because $X^{-1}H$ has a zero diagonal.
The second-order term is $-\tfrac12\tr(X^{-1}HX^{-1}H)=-\tfrac12\cdot2\cdot\tfrac{0.01}{8}=-0.00125$,
which agrees with the exact change to about $8\times10^{-7}$.

\paragraph{Exercise 1.10}
A sketch in numpy.
\begin{verbatim}
import numpy as np
rng = np.random.default_rng(0)
S = np.diag([1., 3.]); mu = np.array([1., 0.])
P = np.array([[2., 1.], [1., 2.]])
d = mu + rng.standard_normal((100000, 2)) @ np.sqrt(S)   # S is diagonal here
print(np.mean(np.einsum('ni,ij,nj->n', d, P, d)))         # close to 10
\end{verbatim}
For a Gaussian $\delta$ the variance of $\delta\T P\delta$ is
$2\tr(P\Sigma P\Sigma)+4\mu\T P\Sigma P\mu=92+28=120$, so the sample mean has
standard error $\sqrt{120/10^5}\approx0.035$, and values within about $0.1$ of
$10$ are expected.
